%% toc: P15 164 164.1
\section{Evidential relevance: ordinal equivalence, direction and the source examples, proofs}
\label{sec:p15-w02c-p15-027-annex}

Use the definitions of Chapter~\ref{sec:p15-w02c-p15-027-relevance} and
Definition~\ref{def:p15-r43-measures}. When one source \(E\) and one
target \(H\) are in view, write \(e=P(E)\), \(h=P(H)\),
\(q=P(H\mid E)\), \(a=P(E\mid H)\) and \(b=P(E\mid\neg H)\).

\emph{Two realization facts} are used repeatedly.
(R1) For \(h\in(0,1)\) and \(q\in[0,1]\), put \(e=\min\{h,1-h\}/2\).
The four numbers \(eq\), \(e(1-q)\), \(h-eq\), \(1-h-e(1-q)\) are
nonnegative, since \(eq\leq e<h\) and \(e(1-q)\leq e<1-h\), and they
sum to one. On a four-point sample space with these masses, let \(E\) be
the first two points and \(H\) the first and third. Then \(P(E)=e\in(0,1)\),
\(P(H)=h\) and \(P(H\mid E)=q\).
(R2) For \(h\in(0,1)\) and \(a,b\in[0,1]\), the masses \(ha\),
\(h(1-a)\), \((1-h)b\), \((1-h)(1-b)\) on four points, with \(H\) the
first two points and \(E\) the first and third, give \(P(H)=h\),
\(P(E\mid H)=a\), \(P(E\mid\neg H)=b\) and \(P(E)=ha+(1-h)b\), which
lies strictly between zero and one unless \(a=b=0\) or \(a=b=1\).

\begin{annexitem}{annex:p15-w02c-p15-027-product}{thm:p15-w02c-p15-027-product}{Definitions 2, 3 and 5}
\begin{proof}
\emph{Item 1.} If \(P(A)>0\), then \(P(B\mid A)=P(AB)/P(A)\), and
multiplying both sides of \(P(B\mid A)\bowtie P(B)\) by the positive
number \(P(A)\) gives \(P(AB)\bowtie P(A)P(B)\); dividing reverses the
step. Multiplication and division by a positive number preserve each of
\(<\), \(=\) and \(>\).

\emph{Item 2.} If \(P(A)=0\), then \(0\leq P(AB)\leq P(A)=0\) by
monotonicity, so \(P(AB)=0=P(A)P(B)\), and the condition of
Definition~3 is the equality case: \(A\) is irrelevant to \(B\).
Definition~2 classifies \(A\) as irrelevant to \(B\) in this case by its
clause \emph{or \(P(A)=0\)}. Together with item~1, the two definitions
give the same classification for every \(P\), \(A\), \(B\).

\emph{Item 3.} \(P(AB)=P(BA)\) and \(P(A)P(B)=P(B)P(A)\), so the
comparison of Definition~3 for \(A\) to \(B\) is literally the comparison
for \(B\) to \(A\). In the three-outcome model, \(P(A)=0\) and
\(P(B)=1/2\). Then \(P(A\mid B)=P(AB)/P(B)=0=P(A)\), so the ratio
comparison makes \(B\) irrelevant to \(A\); and \(P(B\mid A)\) has
denominator zero, so by Definition~1 it is undefined and the ratio
comparison, without the completion clause, assigns no relation of \(A\)
to \(B\).

\emph{Item 4.} The integers \(15,\ldots,44\) are thirty in number; the
multiples of \(3\) among them are \(15,18,\ldots,42\), ten in number, and
the multiples of \(15\) are \(15\) and \(30\). So \(P(B)=1/3\) and
\(P(A)=1/15\). Every multiple of \(15\) is a multiple of \(3\), so
\(AB=A\), \(P(B\mid A)=1\) and \(P(A\mid B)=P(A)/P(B)=1/5\). Hence
\(R_D(A\Yright B)_P=1-1/3=2/3>0\) and
\(R_D(B\Yright A)_P=1/5-1/15=2/15>0\); both directions are positive, by
the sign of D (Theorem~\ref{thm:p15-r43-measures}).

\emph{Item 5.} If \(P(K)=1\), then \(P(X\neg K)\leq P(\neg K)=0\), so
\(P(XK)=P(X)\) and \(P(X\mid K)=P(XK)/P(K)=P(X)\) for every \(X\). The
tautology \(\top\) is true at every outcome, so its event is the whole
space and \(P(\top)=1\). Substituting \(X=AC\), \(X=A\) and \(X=C\) turns
\(P(AC\mid K)\bowtie P(A\mid K)P(C\mid K)\) into
\(P(AC)\bowtie P(A)P(C)\). The equality case is
\((A\perp C\mid K)_P\) on the left and irrelevance in the sense of
Definition~3 on the right; the strict cases are the conditional and
unconditional polarities.

\emph{Item 6.} If \(P(B)>0\), then
\(P(AC\mid B)=P(ABC)/P(B)\) and
\(P(A\mid B)P(C\mid B)=P(AB)P(BC)/P(B)^2\). Multiplying both sides of
the comparison by \(P(B)^2>0\) gives
\(P(ABC)P(B)\bowtie P(AB)P(BC)\), and dividing reverses the step. The
right-hand form uses only unconditional probabilities, so it is defined
for every \(B\). If \(P(B)=0\), then \(P(ABC)\), \(P(AB)\) and \(P(BC)\)
are all at most \(P(B)=0\), so both products are zero.
\end{proof}
\end{annexitem}

\begin{annexitem}{annex:p15-w02c-p15-027-transitivity}{thm:p15-w02c-p15-027-transitivity}{Transitivity}
\begin{proof}
\emph{Item 1.} By Definition~\ref{def:p15-r43-extreme},
\(A\,\mathrm{expos}_P\,B\) gives \(P(A)>0\) and
\(P(B)<1=P(B\mid A)\), so \(P(B\mid A)>P(B)\) and \(A\) is positive to
\(B\). Likewise \(C\,\mathrm{expos}_P\,B\) gives \(P(C)>0\) and
\(P(B\mid C)>P(B)\), so \(C\) is positive to \(B\); multiplying by
\(P(C)\) gives \(P(BC)>P(B)P(C)\), and since this condition is symmetric
in its two events, \(B\) is positive to \(C\). Now
\(A\,\mathrm{exneg}_P\,C\) means \(A\,\mathrm{expos}_P\,\neg C\), that is,
\(P(A)>0\), \(P(\neg C)<1\) and \(P(\neg C\mid A)=1\). The first holds.
The second is \(P(C)>0\). For the third,
\(P(A\neg C)=P(A)-P(AC)=P(A)\), so \(P(\neg C\mid A)=1\). Only
\(P(A)>0\), \(P(C)>0\) and \(P(AC)=0\) were used. Finally
\(P(C\mid A)=P(AC)/P(A)=0<P(C)\), so \(A\) is negative to \(C\).

\emph{Item 2.} Let \(P(A),P(C)>0\), \(P(AC)=0\) and \(B=A\vee C\) with
\(P(B)<1\). Since \(A\) entails \(B\), \(P(AB)=P(A)\), so
\(P(B\mid A)=1>P(B)\) and \(A\,\mathrm{expos}_P\,B\); in the same way
\(C\,\mathrm{expos}_P\,B\). In the three-outcome model,
\(P(A)=P(C)=1/3\), \(A\) and \(C\) are disjoint, and \(P(B)=2/3<1\).

\emph{Item 3.} With \(A\) the first two of four equally weighted outcomes,
\(P(A)=1/2\). If \(B\) is the whole space, \(P(AB)=P(A)\), so
\(P(B\mid A)=1=P(B)\) and \(A\) is not positive to \(B\); it is
irrelevant, since \(P(AB)=P(A)P(B)\). If \(B\) is the first outcome,
\(P(B)=1/4\) and \(P(AB)=1/4\), so \(P(B\mid A)=1/2\), which lies
strictly between \(P(B)\) and one.
\end{proof}
\end{annexitem}

\begin{annexitem}{annex:p15-w02c-p15-027-simpson}{thm:p15-w02c-p15-027-simpson}{A stratified population}
\begin{proof}
The counts are consistent. Among the British, \(32\) were given the drug
and cured, \(40-32=8\) were given the drug only, \(40-32=8\) were cured
only, and \(50-32-8-8=2\) neither. Among the others the four numbers are
\(2\), \(8\), \(8\) and \(50-2-8-8=32\). All are nonnegative. Under
uniform sampling the probability of a proposition about \(\iota\) is the
number of individuals with the property divided by one hundred.

\emph{Item 1.} \(P(B)=50/100\); \(P(A)=(40+10)/100\);
\(P(C)=(40+10)/100\); \(P(AB)=40/100\); \(P(BC)=40/100\);
\(P(AC)=(32+2)/100=17/50\).

\emph{Item 2.} Conditioning on \(B\) divides counts among the British
by fifty: \(P(A\mid B)=P(C\mid B)=4/5\) and
\(P(AC\mid B)=32/50=16/25=(4/5)^2\). Among the others,
\(P(A\mid\neg B)=P(C\mid\neg B)=1/5\) and
\(P(AC\mid\neg B)=2/50=1/25=(1/5)^2\).

\emph{Item 3.} \(\Delta_P(A,B)=2/5-1/4=3/20\) and
\(\Delta_P(B,C)=2/5-1/4=3/20\), both positive.
\(\Delta_P(A,C)=17/50-1/4=9/100>0\). With \(P(B)P(\neg B)=1/4\),
\(\Delta_P(A,B)\Delta_P(B,C)/(P(B)P(\neg B))=(9/400)\cdot4=9/100\),
which is the value of Theorem~\ref{thm:p15-r43-screening} under the
conditional independence of item~2. Finally
\(P(C\mid A)=(17/50)/(1/2)=17/25>1/2\).
\end{proof}
\end{annexitem}

\begin{annexitem}{annex:p15-w02c-p15-027-ranges}{thm:p15-w02c-p15-027-ranges}{Ranges of the seven measures}
\begin{proof}
For each measure we show that every value lies in the stated interval
and that every point of the interval is attained by a model of (R1) or
(R2), in which \(0<P(E),P(H)<1\).

\emph{D.} The completion value is zero. If \(e>0\) and \(q<1\), then
\(R_D=q-h<1\). If \(e>0\) and \(q=1\), then \(h\geq P(EH)=e>0\), so
\(R_D\leq1-e<1\). If \(h=1\), then \(P(E\neg H)=0\) and \(q=1\), so
\(R_D=0\); if \(h<1\), then \(R_D\geq-h>-1\). Hence every value lies in
\((-1,1)\). For \(t\in(-1,1)\), take \(h=(1-t)/2\in(0,1)\) and
\(q=(1+t)/2\in(0,1)\) in (R1); then \(R_D=t\).

\emph{Q.} The completion value is one. Otherwise \(e,h>0\) and
\(R_Q=q/h\) is finite and nonnegative. For \(t\in[0,\infty)\) take
\(h=1/(2+t)\in(0,1/2]\) and \(q=t/(2+t)\in[0,1)\) in (R1); then
\(R_Q=t\). The value \(+\infty\) is never taken.

\emph{L.} The completion value is one. Otherwise \(0<h<1\) and either
\(b>0\), giving \(a/b\in[0,\infty)\), or \(b=0<a\), giving \(+\infty\).
For \(t\in[0,\infty)\) take \(h=1/2\), \(a=t/(1+t)\), \(b=1/(1+t)\) in
(R2); \(P(E)=1/2\) and \(R_L=t\). For \(+\infty\) take \(h=1/2\),
\(a=1\), \(b=0\); \(P(E)=1/2\).

\emph{QL and LL.} The natural logarithm, with \(\ln0=-\infty\) and
\(\ln(+\infty)=+\infty\), is a strictly increasing bijection of
\([0,\infty)\) onto \([-\infty,\infty)\) and of \([0,\infty]\) onto
\([-\infty,\infty]\). The completions give \(\ln1=0\). The ranges of QL
and LL are therefore the images of the ranges of Q and L, attained by the
same models.

\emph{KO.} The completion value is zero. Otherwise \(0<e,h<1\), so
\(e=ha+(1-h)b>0\) forces \(a+b>0\), and \(|a-b|\leq a+b\) gives
\(R_{KO}\in[-1,1]\). For \(t\in[-1,1]\) take \(h=1/2\), \(a=(1+t)/2\),
\(b=(1-t)/2\) in (R2); \(P(E)=1/2\) and \(R_{KO}=t\).

\emph{Z.} The completion value is zero. If \(q>h\), then
\(0<(q-h)/(1-h)\leq(1-h)/(1-h)=1\). If \(q\leq h\), then
\(0\geq(q-h)/h\geq-h/h=-1\). For \(t\in[-1,1]\) take \(h=1/2\) and
\(q=(1+t)/2\) in (R1). If \(t>0\), then \(q>h\) and
\(R_Z=(t/2)/(1/2)=t\); if \(t\leq0\), then \(q\leq h\) and again
\(R_Z=(t/2)/(1/2)=t\).

Each of the seven sets is an interval of \(\mathbb R^*\), and every
witness above uses a four-point space with \(0<P(E),P(H)<1\).

\emph{The polarity indicator.} \(R_D\) is a function of the joint table:
when \(e>0\), both \(q=P(EH)/e\) and \(h\) are read from the table, and
\(e=0\) is itself visible in the table. Hence so is
\(\operatorname{sgn}R_D\). When \(e>0\),
\(R_D=q-h=\Delta_P(E,H)/e\) has the sign of the covariance; when
\(e=0\), the covariance is zero and the value is zero. So
\(\operatorname{sgn}R_D\) exceeds, equals or falls below zero exactly
when \(E\) is positive, irrelevant or negative to \(H\): it satisfies
Quality Representation with \(\iota=0\). By the D case, \(R_D\) takes
the values \(-1/2,0,1/2\), so \(\operatorname{sgn}R_D\) takes all of
\(-1,0,1\) and nothing else. The set \(\{-1,0,1\}\) omits \(1/2\),
which lies between two of its members, so it is not an interval.
\end{proof}
\end{annexitem}

\begin{annexitem}{annex:p15-w02c-p15-027-ordinal}{thm:p15-w02c-p15-027-ordinal}{Ordinal equivalence and increasing transformation}
\begin{proof}
\emph{A joint realization.} Let \(\pi,\pi'\in\Sigma\). On the sixteen-point
space \(\{1,2,3,4\}^2\) give the point \((i,j)\) the mass
\(\pi_i\pi'_j\). Let \(A_1\) be the points with \(i\in\{1,2\}\), \(B_1\)
those with \(i\in\{1,3\}\), \(A_2\) those with \(j\in\{1,2\}\) and \(B_2\)
those with \(j\in\{1,3\}\). Summing the product masses over \(j\) gives
\(\pi_P(A_1,B_1)=\pi\), and over \(i\) gives \(\pi_P(A_2,B_2)=\pi'\).
So any two joint tables occur together under one probability measure.

\emph{Item 1.} Fix \(P\) and two pairs, and let
\(x=R_T(A\Yright B)_P\), \(y=R_T(C\Yright D)_P\), both in
\(\operatorname{ran}T\). Because \(f\) is strictly increasing there,
\(x<y\) implies \(f(x)<f(y)\), \(x=y\) implies \(f(x)=f(y)\), and
\(x>y\) implies \(f(x)>f(y)\). Exactly one of the three relations holds
between \(x\) and \(y\), and exactly one between \(f(x)\) and \(f(y)\),
so for each \(\bowtie\) the relation holds between \(x,y\) exactly when
it holds between \(f(x),f(y)\). These are the values of \(T'\).

\emph{Item 2.} By hypothesis \(R_T=\phi_T\circ\pi_P\) and
\(R_{T'}=\phi_{T'}\circ\pi_P\) for functions \(\phi_T,\phi_{T'}\) on
\(\Sigma\). Let \(t\in\operatorname{ran}T\) and choose a table \(\pi\) with
\(\phi_T(\pi)=t\); put \(f(t)=\phi_{T'}(\pi)\). If \(\pi'\) is another
table with \(\phi_T(\pi')=t\), realize \(\pi\) and \(\pi'\) together as
above. The two pairs then have equal \(T\) values under one measure, so
ordinal equivalence with \(\bowtie\) taken as equality gives
\(\phi_{T'}(\pi)=\phi_{T'}(\pi')\). Hence \(f\) is well defined, and
\(R_{T'}=f(R_T)\) holds for every \(P,A,B\) because both sides are
computed from the same table. If \(t<u\) in \(\operatorname{ran}T\),
choose tables \(\pi,\pi'\) with these \(T\) values, realize them together
and apply ordinal equivalence with \(<\): \(f(t)<f(u)\). Every value of
\(T'\) is \(\phi_{T'}(\pi)=f(\phi_T(\pi))\) for some \(\pi\), so \(f\) is
onto \(\operatorname{ran}T'\). Any \(g\) with \(R_{T'}=g(R_T)\) agrees
with \(f\) at every attained value, which is every point of its domain.

\emph{Item 3.} Let \(I=\operatorname{ran}T\) and
\(J=\operatorname{ran}T'\) be intervals, and let \(x\in I\) not be the
least point of \(I\). Put \(s=\sup\{f(y):y\in I,\ y<x\}\leq f(x)\), and
suppose \(s<f(x)\). Since \(\mathbb R^*\) is densely ordered, there is
\(w\) with \(s<w<f(x)\). Some \(y_0<x\) lies in \(I\), and
\(f(y_0)\leq s<w<f(x)\), so \(w\in J\) because \(J\) is an interval. As
\(f\) is onto \(J\), \(w=f(z)\) for some \(z\in I\). If \(z<x\), then
\(f(z)\leq s<w\); if \(z\geq x\), then \(f(z)\geq f(x)>w\). Both are
impossible, so \(s=f(x)\). In the same way,
\(\inf\{f(y):y\in I,\ y>x\}=f(x)\) whenever \(x\) is not the greatest
point of \(I\). Now let \(u<f(x)<v\). If \(x\) is not least, then
\(s=f(x)>u\) gives \(y_1<x\) in \(I\) with \(f(y_1)>u\); if \(x\) is
least, put no lower restriction. Symmetrically obtain \(y_2>x\) with
\(f(y_2)<v\), or no upper restriction. Let \(N\) be the set of
\(y\in I\) with \(y_1<y<y_2\), omitting a bound that was not chosen;
it is a neighbourhood of \(x\) in \(I\). For \(y\in N\), monotonicity
gives \(f(y)>f(y_1)>u\) when \(y_1\) was chosen and
\(f(y)\geq f(x)>u\) when \(x\) is least; symmetrically \(f(y)<v\).
So the preimage of each open interval around \(f(x)\) contains
a neighbourhood of \(x\) in \(I\): \(f\) is continuous at \(x\).

\emph{Item 4.} For a fixed \(P\), \(c(P)\) is a positive constant, and
multiplying by a positive constant preserves each of \(<,=,>\). So the
comparisons of Definition~\ref{def:p15-w02c-p15-027-ordinal}, made
within one \(P\), agree for \(D\) and \(T'\), and \(T'\) has the sign of
\(R_D\), which gives Quality Representation with neutral value zero
(\(R_D\) satisfies it, as shown for its sign in
Annex~\ref{annex:p15-w02c-p15-027-ranges}). The models of the D case of
that annex use four points and give every value in \((-1,1)\) with
\(c=1\). Adding a fifth point of mass zero leaves every table unchanged
and gives the same D values with \(c=2\). So \(T'\) attains all of
\((-1,1)\cup(-2,2)=(-2,2)\), and \(|c\,R_D|<2\) always, so its range is
the interval \((-2,2)\). Finally, the four-point and five-point models
for \(R_D=1/2\) give \(T'\) the values \(1/2\) and \(1\). A function
\(f\) with \(R_{T'}=f(R_D)\) would need \(f(1/2)\) to equal both. (The
same pair shows that \(T'\) is not a function of the joint table.)

\emph{Item 5.} The extended logarithm is strictly increasing on
\([0,\infty]\), which contains \(\operatorname{ran}L\), and
\(R_{LL}=\ln R_L\), so item~1 gives \(L\sim_{\uparrow}LL\). By
Theorem~\ref{thm:p15-w02c-p15-027-ranges}, \(R_L\) attains \(+\infty\);
a function \(f:\mathbb R\to\mathbb R\) is not defined there, so it cannot
give \(R_{LL}=f(R_L)\). Suppose \(f:\mathbb R^*\to\mathbb R^*\) is
strictly increasing with \(R_{LL}=f(R_L)\). \(R_L\) also attains zero,
where \(R_{LL}=-\infty\), so \(f(0)=-\infty\); then \(f(-1)<f(0)=-\infty\),
which no member of \(\mathbb R^*\) satisfies.
\end{proof}
\end{annexitem}

\begin{annexitem}{annex:p15-w02c-p15-027-classes}{thm:p15-w02c-p15-027-classes}{Members of the class of D}
\begin{proof}
\emph{Item 1.} POP-D and MDHD each have the form: for the stated
events, \(R(X)\bowtie R(Y)\) exactly when a probability condition
\(u\bowtie v\) holds, for every \(\bowtie\) among \(<,=,>\). D satisfies
both on the comparison domain (Theorem~\ref{thm:p15-r43-characterization}).
If \(f\) is strictly increasing on \((-1,1)\), which contains every value
of \(R_D\), then \(f(R_D(X))\bowtie f(R_D(Y))\) exactly when
\(R_D(X)\bowtie R_D(Y)\): a strictly increasing function preserves and
reflects \(<\), and it is injective, so it preserves and reflects \(=\).
So \(f\circ R_D\) satisfies both axioms. The same equivalence, applied to
any two values of \(R_D\) under one \(P\), gives
\(f\circ R_D\sim_{\uparrow}D\) by item~1 of
Theorem~\ref{thm:p15-w02c-p15-027-ordinal}.

\emph{Item 2.} Write the joint table of \(E\) and \(H\) as
\[
 (x,y,z,w)=\bigl(P(EH),\,P(E\neg H),\,P(\neg EH),\,P(\neg E\neg H)\bigr).
\]
Where
\(x+y>0\),
\(R_D(E\Yright H)_P=x/(x+y)-(x+z)\), a rational function with nonzero
denominator, hence differentiable; in particular it is differentiable
where all four cells are positive, and so is its cube, as a product of
differentiable functions. The cube \(t\mapsto t^3\) is strictly
increasing and continuous and maps \((-1,1)\) onto \((-1,1)\); since
\(R_D\) attains every value of \((-1,1)\) with source and target
probabilities strictly between zero and one
(Theorem~\ref{thm:p15-w02c-p15-027-ranges}), \(R_{D^3}\) attains every
value of \((-1,1)\) and no other. Its neutral value is \(0^3=0\), and
\(t^3\) has the sign of \(t\), so it represents polarity exactly as D
does. By item~1 it satisfies POP-D and MDHD. Fact (R1) with
\(h=1/4\) and \(q=3/4\) gives a model with \(0<P(E)<1\),
\(P(H)=1/4\) and \(P(H\mid E)=3/4\), where \(R_D=1/2\) and
\(R_{D^3}=1/8\).

\emph{Item 3.} The map \(g(t)=t\) for \(t\geq0\), \(g(t)=2t\) for
\(t<0\), is strictly increasing on \((-1,1)\), continuous, has
\(g(0)=0\) and the sign of \(t\), and maps \((-1,1)\) onto \((-2,1)\).
By item~1, \(R_{D^\sharp}=g\circ R_D\) satisfies POP-D and MDHD; it
represents polarity with neutral value zero; and, as in item~2, its
range is \(g((-1,1))=(-2,1)\), an interval. At the table
\((1/4,1/4,1/4,1/4)\), \(\Delta=1/4-(1/2)(1/2)=0\). Move mass \(t\) from
\(y\) to \(x\), giving \((1/4+t,1/4-t,1/4,1/4)\) for \(|t|<1/4\), which
keeps every cell positive and \(x+y=1/2\). Along this path
\(R_D=(1/4+t)/(1/2)-(1/2+t)=t\), so \(R_{D^\sharp}=t\) for \(t\geq0\)
and \(2t\) for \(t<0\). The one-sided derivatives at \(t=0\) are \(1\)
and \(2\). A function differentiable at the table would have a
derivative along this path equal to a single linear functional applied
to the direction \((1,-1,0,0)\), so the two one-sided derivatives would
coincide. Hence \(R_{D^\sharp}\) is not differentiable there.
\end{proof}
\end{annexitem}

\begin{annexitem}{annex:p15-w02c-p15-027-direction}{thm:p15-w02c-p15-027-direction}{Direction formulas and the three-way model}
\begin{proof}
Let \(x=P(A)\), \(y=P(B)\), both in \((0,1)\), and
\(\delta=P(AB)-xy\ne0\).

\emph{Formulas.} \(R_D(A\Yright B)_P=P(AB)/x-y=\delta/x\).
Next \(P(B\mid\neg A)=(y-P(AB))/(1-x)\), so
\(P(B\mid\neg A)-y=(y-P(AB)-y+xy)/(1-x)=-\delta/(1-x)\), and therefore
\[
 \delta p(A\Yright B)=\bigl(P(B\mid A)-y\bigr)-\bigl(P(B\mid\neg A)-y\bigr)
 =\frac{\delta}{x}+\frac{\delta}{1-x}=\frac{\delta}{x(1-x)}.
\]
\(R_Q(A\Yright B)_P=P(AB)/(xy)=1+\delta/(xy)\). For Z,
\(q-h=\delta/x\); if \(\delta>0\), then \(q>h\) and
\(R_Z=(\delta/x)/(1-y)\); if \(\delta<0\), then \(q<h\) and
\(R_Z=(\delta/x)/y\). The formulas for \(B\Yright A\) follow by
exchanging \(x\) and \(y\), since \(\delta\) is symmetric.

\emph{Comparisons.} The expressions for Q, and for Z when \(\delta<0\),
are symmetric in \(x,y\), so the two directions have equal values. For
D the absolute values are \(|\delta|/x\) and \(|\delta|/y\), and the
first is larger exactly when \(x<y\). For Z with \(\delta>0\) the
denominators are \(x(1-y)\) and \(y(1-x)\), whose difference is
\(x-y\); the value for \(A\Yright B\) is larger exactly when \(x<y\).
For \(\delta p\) the denominators are \(x(1-x)\) and \(y(1-y)\), and
\(z(1-z)=\tfrac14-(z-\tfrac12)^2\), so the absolute value for
\(A\Yright B\) is larger exactly when \(|x-\tfrac12|>|y-\tfrac12|\).

\emph{The model.} The four masses are nonnegative and sum to
\((17+3+28+2)/50=1\). Then \(x=20/50=2/5\), \(y=45/50=9/10\),
\(P(AB)=17/50\) and \(\delta=17/50-18/50=-1/50<0\), so \(A\) is negative
to \(B\). The formulas give
\(R_D(A\Yright B)=(-1/50)/(2/5)=-1/20\) and
\(R_D(B\Yright A)=(-1/50)/(9/10)=-1/45\);
\(\delta p(A\Yright B)=(-1/50)/(6/25)=-1/12\) and
\(\delta p(B\Yright A)=(-1/50)/(9/100)=-2/9\);
\(R_Q=1-(1/50)/(9/25)=17/18\) in both directions; and
\(R_Z=(-1/50)/(9/25)=-1/18\) in both directions. For L,
\(P(A\mid B)=(17/50)/(9/10)=17/45\) and
\(P(A\mid\neg B)=(3/50)/(1/10)=3/5\), so
\(R_L(A\Yright B)=(17/45)/(3/5)=17/27\); and
\(P(B\mid A)=(17/50)/(2/5)=17/20\),
\(P(B\mid\neg A)=(28/50)/(3/5)=14/15\), so
\(R_L(B\Yright A)=(17/20)/(14/15)=51/56\). Distances from the neutral
values: D gives \(1/20>1/45\); \(\delta p\) gives \(1/12<2/9\); Q and Z
give equality. For L, \(27\cdot51=1377>952=56\cdot17\), so
\(27/17>56/51\) and \(|\ln(17/27)|>|\ln(51/56)|\).
\end{proof}
\end{annexitem}

\begin{annexitem}{annex:p15-w02c-p15-027-likewise}{thm:p15-w02c-p15-027-likewise}{L in the two directions}
\begin{proof}
\emph{Preliminaries.} If \(u\in\{0,1\}\), then \(x=0=uv\) when \(u=0\)
and \(x=v=uv\) when \(u=1\), since \(P(\neg A B)\leq P(\neg A)=0\); the
same holds for \(v\). Either way \(A\) and \(B\) would be independent,
so \(0<u,v<1\), and \(x\ne uv\) by hypothesis. Always
\(0\leq x\leq\min\{u,v\}\), and \(x=u+v-P(A\vee B)\geq u+v-1\). Since
\(\min\{u,v\}>0\), the four cases \(0<x<\min\{u,v\}\);
\(x=\min\{u,v\}\) with \(u<v\), with \(v<u\), or with \(u=v\); and
\(x=0\) are mutually exclusive and exhaust the possibilities. In each
case both targets are nondegenerate, so
Definition~\ref{def:p15-r43-measures} computes \(R_L(A\Yright B)_P\)
from the likelihoods
\[
 P(A\mid B)=\frac{x}{v},\qquad P(A\mid\neg B)=\frac{u-x}{1-v},
\]
and \(R_L(B\Yright A)_P\) from
\[
 P(B\mid A)=\frac{x}{u},\qquad P(B\mid\neg A)=\frac{v-x}{1-u}.
\]

\emph{Case 1.} All four likelihoods are positive, so
\[
 R_L(A\Yright B)_P=\frac{x/v}{(u-x)/(1-v)}=\frac{x(1-v)}{v(u-x)},\qquad
 R_L(B\Yright A)_P=\frac{x(1-u)}{u(v-x)},
\]
both finite and positive. Over the common denominator
\(uv(u-x)(v-x)>0\) the difference has numerator
\[
 x\bigl[u(1-v)(v-x)-v(1-u)(u-x)\bigr].
\]
Expanding, \(u(1-v)(v-x)=uv-ux-uv^2+uvx\) and
\(v(1-u)(u-x)=uv-vx-u^2v+uvx\), so the bracket equals
\(x(v-u)+uv(u-v)=(v-u)(x-uv)\), and the numerator is
\(x(v-u)(x-uv)\), which is the displayed difference. Further,
\[
\begin{aligned}
 R_L(A\Yright B)_P-1&=\frac{x(1-v)-v(u-x)}{v(u-x)}=\frac{x-uv}{v(u-x)},\\
 R_L(B\Yright A)_P-1&=\frac{x-uv}{u(v-x)},
\end{aligned}
\]
so both values exceed one when \(x>uv\) and both lie in \((0,1)\) when
\(x<uv\). If \(x>uv\), then \(|\ln r|=\ln r\) is strictly increasing in
\(r>1\), and the difference has the sign of \(v-u\); so
\(|\ln R_L(A\Yright B)_P|\) is greater than, equal to or less than
\(|\ln R_L(B\Yright A)_P|\) according as \(u<v\), \(u=v\) or \(u>v\).
If \(x<uv\), then \(|\ln r|=-\ln r\) is strictly decreasing in
\(r\in(0,1)\), and the difference has the sign of \(-(v-u)=u-v\); so
\(R_L(A\Yright B)_P\) is the smaller value, and has the larger distance,
exactly when \(u<v\), with equality exactly when \(u=v\).

\emph{Case 2.} Let \(x=u<v\). Then \(P(A)=u>0\), \(P(B)=v<1\) and
\(P(B\mid A)=x/u=1\), so \(A\,\mathrm{expos}_P\,B\) by
Definition~\ref{def:p15-r43-extreme}. The likelihoods for
\(A\Yright B\) are \(P(A\mid B)=u/v>0\) and
\(P(A\mid\neg B)=(u-u)/(1-v)=0\), so the definition gives
\(R_L(A\Yright B)_P=+\infty\). The likelihoods for \(B\Yright A\) are
\(P(B\mid A)=1\) and \(P(B\mid\neg A)=(v-u)/(1-u)\), which is positive
because \(v>u\) and less than one because \(v<1\). Hence
\[
 R_L(B\Yright A)_P=\frac{1}{(v-u)/(1-u)}=\frac{1-u}{v-u},
\]
finite, and greater than one since \(1-u>v-u\). The distances are
\(+\infty\) and \(\ln((1-u)/(v-u))<+\infty\), so the source of smaller
probability, \(A\), has the larger distance, as the final clause
requires for \(u<v\). The mirror case \(x=v<u\) is the same computation
with \(A\) and \(B\) exchanged: \(B\,\mathrm{expos}_P\,A\),
\(R_L(B\Yright A)_P=+\infty\), and
\(R_L(A\Yright B)_P=1/P(A\mid\neg B)=(1-v)/(u-v)\in(1,+\infty)\);
the larger distance again belongs to the less probable source, here
\(B\).

\emph{Case 3.} Let \(x=u=v\). Then \(P(A\mid\neg B)=0<u/v=1=P(A\mid B)\)
and \(P(B\mid\neg A)=0<1=P(B\mid A)\), so both values are \(+\infty\),
the distances are equal, and \(u=v\).

\emph{Case 4.} Let \(x=0\). Then \(P(\neg B\mid A)=1\) with
\(P(A)>0\) and \(P(\neg B)<1\), which is \(A\,\mathrm{exneg}_P\,B\), and
symmetrically \(B\,\mathrm{exneg}_P\,A\). The likelihoods \(P(A\mid B)\)
and \(P(B\mid A)\) are zero, while \(P(A\mid\neg B)=u/(1-v)\) and
\(P(B\mid\neg A)=v/(1-u)\) are positive, so both L values are
\(0\) and both distances are \(+\infty\), whatever \(u\) and \(v\) are.
For D, \(R_D(A\Yright B)_P=P(B\mid A)-v=-v\) and
\(R_D(B\Yright A)_P=-u\), with absolute values \(v\) and \(u\).

\emph{The numerical instance after the theorem.} With \(u=1/4\),
\(v=1/2\), \(x=1/4\) we are in Case~2, and
\((1-u)/(v-u)=(3/4)/(1/4)=3\). With \(x=0\), \(x\geq u+v-1=-1/4\) holds
and \(x\ne uv=1/8\), so we are in Case~4, and
\(R_D(A\Yright B)_P=-1/2\), \(R_D(B\Yright A)_P=-1/4\).
\end{proof}
\end{annexitem}

\begin{annexitem}{annex:p15-w02c-p15-027-covariance}{thm:p15-w02c-p15-027-covariance}{Carnap's measure}
\begin{proof}
\emph{Item 1.} \(R_C(E\Yright H)_P=P(EH)-P(E)P(H)=\Delta_P(E,H)\), and
\(\Delta_P(E,H)=\Delta_P(H,E)\) because conjunction and multiplication
are commutative; this is Definition~7 for \(R_C\). With
\((x,y,z,w)\) the cells \(EH\), \(E\neg H\), \(\neg EH\), \(\neg E\neg H\),
\(P(E)=x+y\), \(P(H)=x+z\) and \(x+y+z+w=1\), so
\[
 x-(x+y)(x+z)=x(1-x-y-z)-yz=xw-yz .
\]
This is a polynomial in the cells, hence differentiable. Since
\(x+w\leq1\), \(xw\leq((x+w)/2)^2\leq1/4\), so \(R_C\leq xw\leq1/4\);
likewise \(yz\leq1/4\), so \(R_C\geq-yz\geq-1/4\). For
\(s\in[0,1/2]\) the table \((s,1/2-s,1/2-s,s)\) gives
\(R_C=s^2-(1/2-s)^2=s-1/4\), which runs through \([-1/4,1/4]\); so the
range is exactly that interval. The sign of \(\Delta_P(E,H)\) is the
polarity of \(E\) to \(H\) by Definition~3, so Quality Representation
holds with neutral value zero, and \(R_C\) satisfies both conditions of
Definition~\ref{def:p15-w02c-p15-027-measure}.

\emph{Item 2.} Each outcome has mass \(1/40\). \(E\) has twenty
outcomes, \(H_1\) four, all in \(E\), and \(H_2\) twenty, of which
fifteen (outcomes \(5\)--\(19\)) are in \(E\). So \(P(E)=1/2\),
\(P(H_1)=1/10\), \(P(EH_1)=1/10\), \(P(H_2)=1/2\), \(P(EH_2)=3/8\),
\(P(E\mid H_1)=1\) and \(P(E\mid H_2)=(3/8)/(1/2)=3/4\). Then
\(\Delta_P(E,H_1)=1/10-1/20=1/20\) and
\(\Delta_P(E,H_2)=3/8-1/4=1/8\), and item~1 gives the equal values in the
other direction. LIK-D for source \(E\) requires
\(R(E\Yright H_1)_P>R(E\Yright H_2)_P\) because
\(P(E\mid H_1)>P(E\mid H_2)\); \(R_C\) gives the opposite strict
inequality. POP-D for target \(E\) requires
\(R(H_1\Yright E)_P>R(H_2\Yright E)_P\) because
\(P(E\mid H_1)>P(E\mid H_2)\); again \(R_C\) gives the opposite. All four
events have probability strictly between zero and one, so both
comparisons lie in the comparison domain.

\emph{Item 3.} \(R_Q(H_i\Yright E)_P=P(E\mid H_i)/P(E)\), which is
\(1/(1/2)=2\) for \(i=1\) and \((3/4)/(1/2)=3/2\) for \(i=2\). Under one
\(P\), Q puts \(H_1\Yright E\) above \(H_2\Yright E\) and \(R_C\) puts it
below, so the condition of
Definition~\ref{def:p15-w02c-p15-027-ordinal} fails for this pair of
comparisons and the two measures are not ordinally equivalent. Q is
amount-symmetric (Theorem~\ref{thm:p15-r43-measures}).
\end{proof}
\end{annexitem}

\begin{annexitem}{annex:p15-w02c-p15-027-negation}{thm:p15-w02c-p15-027-negation}{Negation in the two positions}
\begin{proof}
Write \(e=P(E)\), \(h=P(H)\). The neutral values are \(\iota_D=\iota_Z=0\)
and \(\iota_L=1\), and \(\mathrm{inv}_T(\iota_T)=\iota_T\) in each case.

\emph{Item 1, D.} If \(e=0\), both \(R_D(E\Yright H)_P\) and
\(R_D(E\Yright\neg H)_P\) are the completion \(0\), and \(-0=0\). If
\(e>0\), \(P(\neg H\mid E)=1-q\) and \(P(\neg H)=1-h\), so
\(R_D(E\Yright\neg H)_P=(1-q)-(1-h)=-(q-h)=-R_D(E\Yright H)_P\).

\emph{Item 1, L.} \(\neg H\) is certain or impossible exactly when \(H\)
is, and then both values are the completion \(1=1/1\). Let \(0<h<1\),
\(a=P(E\mid H)\), \(b=P(E\mid\neg H)\). The likelihoods of \(E\) for
the target \(\neg H\) are \(P(E\mid\neg H)=b\) and
\(P(E\mid\neg\neg H)=a\), so the pair is exchanged. If \(a=b=0\), both
values are the completion \(1\). If \(a=0<b\), then
\(R_L(E\Yright H)_P=0\) and \(R_L(E\Yright\neg H)_P=+\infty=1/0\). If
\(b=0<a\), then \(R_L(E\Yright H)_P=+\infty\) and
\(R_L(E\Yright\neg H)_P=0=1/(+\infty)\). If \(a,b>0\), the values are
\(a/b\) and \(b/a\).

\emph{Item 1, Z.} If \(e=0\) or \(h\in\{0,1\}\), the same degeneracy
holds for \(\neg H\), and both values are the completion \(0\). Let
\(e>0\), \(0<h<1\), \(\bar q=1-q\), \(\bar h=1-h\). If \(q>h\), then
\(\bar q<\bar h\), and the second branch for \(\neg H\) gives
\((\bar q-\bar h)/\bar h=-(q-h)/(1-h)=-R_Z(E\Yright H)_P\). If \(q<h\),
then \(\bar q>\bar h\), and the first branch gives
\((\bar q-\bar h)/(1-\bar h)=-(q-h)/h=-R_Z(E\Yright H)_P\). If \(q=h\),
then \(\bar q=\bar h\) and both values are \(0\).

\emph{Item 2.} Let \(e\in\{0,1\}\). Then \(P(\neg E)=1-e\in\{0,1\}\)
too. A source of probability zero gets the completion \(\iota_T\) for D
and Z, and for L both likelihoods vanish, giving \(1\) (or the target
completion \(1\)). A source of probability one has
\(P(H\mid E)=P(EH)/1=h\), since \(P(\neg E H)\leq P(\neg E)=0\); so
\(R_D=0\), \(R_Z=0\) (the second branch with \(q=h\), or the completion
when \(h\in\{0,1\}\)), and for L, when \(0<h<1\), both likelihoods
equal one, giving \(1\). Let \(h\in\{0,1\}\) instead. For any source
\(X\) of positive probability, \(P(H\mid X)=h\) (as in the proof of
Theorem~\ref{thm:p15-w02c-p15-027-mixture}, item~1), so
\(R_D(X\Yright H)_P=0\); and the completions give \(R_L=1\), \(R_Z=0\).
A source of probability zero gets the neutral value as before. In all
these cases both \(R_T(E\Yright H)_P\) and \(R_T(\neg E\Yright H)_P\)
are \(\iota_T\), and \(\iota_T=\mathrm{inv}_T(\iota_T)\).

\emph{Item 3, preliminaries.} Let \(0<e,h<1\) and \(d=q-h\). Then
\(P(\neg E)=1-e>0\) and
\[
 q'=P(H\mid\neg E)=\frac{h-eq}{1-e},\qquad
 q'-h=\frac{h-eq-h+eh}{1-e}=-\frac{e}{1-e}\,d .
 \tag{\(\dagger\)}
\]

\emph{Item 3, D.} By \((\dagger)\),
\(R_D(\neg E\Yright H)_P=-\frac{e}{1-e}R_D(E\Yright H)_P\). This equals
\(-R_D(E\Yright H)_P\) if and only if \(d=0\) or \(e/(1-e)=1\), that is,
\(q=h\) or \(e=1/2\).

\emph{Item 3, L.} The likelihoods of \(\neg E\) are \(1-a\) and \(1-b\).
If \(a=b\), then \(E\) is irrelevant to \(H\) and so is \(\neg E\)
(Theorem~\ref{thm:p15-r43-polarity}); \(a=b\in\{0,1\}\) would force
\(e\in\{0,1\}\), so \(a=b\in(0,1)\) and both values are \(1\), and
inversion holds. Let \(a\ne b\); we show that inversion holds if and
only if \(a+b=1\).
If \(b=0\), then \(a>0\), \(R_L(E\Yright H)_P=+\infty\) with inverse
\(0\), and \(R_L(\neg E\Yright H)_P=(1-a)/1=1-a\), which is \(0\) iff
\(a=1\), that is, iff \(a+b=1\).
If \(a=0\), then \(b>0\), \(R_L(E\Yright H)_P=0\) with inverse
\(+\infty\), and \(R_L(\neg E\Yright H)_P=1/(1-b)\) when \(b<1\) and
\(+\infty\) when \(b=1\); it is \(+\infty\) iff \(b=1\), that is, iff
\(a+b=1\).
If \(b=1\) and \(0<a<1\), then \(R_L(\neg E\Yright H)_P=+\infty\)
(likelihoods \(1-a>0\) and \(0\)), while
\(1/R_L(E\Yright H)_P=1/a\) is finite; inversion fails, and
\(a+b\ne1\).
If \(a=1\) and \(0<b<1\), then \(R_L(\neg E\Yright H)_P=0\), while
\(1/R_L(E\Yright H)_P=b>0\); inversion fails, and \(a+b\ne1\).
Otherwise \(a,b\in(0,1)\), all four likelihoods are positive, and
\((1-a)/(1-b)=b/a\) iff \(a-a^2=b-b^2\) iff \((a-b)(1-a-b)=0\), which
for \(a\ne b\) is \(a+b=1\).

\emph{Item 3, Z.} If \(d=0\), then \(q'=h\) by \((\dagger)\), and both
values are \(0\). Let \(d>0\). Then \(R_Z(E\Yright H)_P=d/(1-h)\), and
by \((\dagger)\) \(q'<h\), so
\(R_Z(\neg E\Yright H)_P=(q'-h)/h=-ed/((1-e)h)\). This equals
\(-d/(1-h)\) iff \(e(1-h)=(1-e)h\) iff \(e=h\). Let \(d<0\). Then
\(R_Z(E\Yright H)_P=d/h\), and \(q'>h\), so
\(R_Z(\neg E\Yright H)_P=(q'-h)/(1-h)=-ed/((1-e)(1-h))\). This equals
\(-d/h\) iff \(eh=(1-e)(1-h)\) iff \(e+h=1\).

\emph{Item 4.} First model: the masses are nonnegative and sum to
\((3+1+5+7)/16=1\); \(h=(3+5)/16=1/2\), \(e=(3+1)/16=1/4\),
\(q=(3/16)/(1/4)=3/4\), and \(q'=(5/16)/(3/4)=5/12\). Since \(q>h\),
\(R_Z(E\Yright H)_P=(3/4-1/2)/(1/2)=1/2\). For \(\neg H\),
\(P(\neg H\mid E)=1/4<1/2=P(\neg H)\), so
\(R_Z(E\Yright\neg H)_P=(1/4-1/2)/(1/2)=-1/2\). Since \(q'<h\),
\(R_Z(\neg E\Yright H)_P=(5/12-1/2)/(1/2)=-1/6\). For D,
\(R_D(E\Yright H)_P=1/4\) and \(R_D(\neg E\Yright H)_P=5/12-1/2=-1/12\).
For L, \(a=(3/16)/(1/2)=3/8\) and \(b=(1/16)/(1/2)=1/8\), so
\(R_L(E\Yright H)_P=3\), with inverse \(1/3\), and
\(R_L(\neg E\Yright H)_P=(5/8)/(7/8)=5/7\). These agree with item~3:
\(q>h\) and \(e=1/4\ne1/2=h\) for Z; \(e\ne1/2\) for D;
\(a+b=1/2\ne1\) for L.

Second model, masses \(3/16\), \(5/16\), \(1/16\) and \(7/16\) for
\(EH\), \(E\neg H\), \(\neg EH\) and \(\neg E\neg H\). Then
\(e=8/16=1/2\), \(h=4/16=1/4\),
\(q=(3/16)/(1/2)=3/8\), \(q'=(1/16)/(1/2)=1/8\). D:
\(3/8-1/4=1/8\) and \(1/8-1/4=-1/8\). Z: \(q>h\), so
\((1/8)/(3/4)=1/6\); \(q'<h\), so \((-1/8)/(1/4)=-1/2\ne-1/6\). L:
\(a=(3/16)/(1/4)=3/4\), \(b=(5/16)/(3/4)=5/12\), so
\(R_L(E\Yright H)_P=9/5\) and
\(R_L(\neg E\Yright H)_P=(1/4)/(7/12)=3/7\ne5/9\). In item~3 terms,
\(e=1/2\) for D, while \(e\ne h\) for Z and \(a+b=7/6\) for L.

Third model, masses \(1/8,1/8,1/8,5/8\): \(e=h=1/4\), \(q=1/2\),
\(q'=(1/8)/(3/4)=1/6\). Z: \((1/2-1/4)/(3/4)=1/3\) and
\((1/6-1/4)/(1/4)=-1/3\). D: \(1/4\) and \(1/6-1/4=-1/12\). L:
\(a=1/2\), \(b=1/6\), so \(R_L(E\Yright H)_P=3\), with inverse \(1/3\),
and \(R_L(\neg E\Yright H)_P=(1/2)/(5/6)=3/5\). In item~3 terms,
\(q>h\) and \(e=h\) for Z, \(e\ne1/2\) for D, \(a+b=2/3\) for L.
\end{proof}
\end{annexitem}

\begin{annexitem}{annex:p15-w02c-p15-027-equivalents}{thm:p15-w02c-p15-027-equivalents}{Equivalent targets and the raven values}
\begin{proof}
If \(G\) and \(H\) are logically equivalent, they have the same truth
event, so \(G\mathbin{\triangle}H=\varnothing\). Suppose
\(P(G\mathbin{\triangle}H)=0\). For any event \(X\), the events \(XG\)
and \(XH\) differ only inside \(G\mathbin{\triangle}H\), so
\(P(XG)=P(XH)\); and \(\neg G\mathbin{\triangle}\neg H=G\mathbin{\triangle}H\),
so likewise \(P(X\neg G)=P(X\neg H)\). Taking \(X=B\) and \(X=\neg B\)
gives \(\pi_P(B,G)=\pi_P(B,H)\), and a function of the joint table takes
the same value on both.

In the raven model, \(P(H)=1/2\). Given \(H\), the cells raven-black,
nonraven-nonblack and nonraven-black have masses \(1/100\), \(4/5\) and
\(19/100\), with raven-nonblack zero. Given \(\neg H\), raven-nonblack
and raven-black have \(1/200\) each, nonraven-nonblack \(159/200\) and
nonraven-black \(39/200\). All are nonnegative and each set sums to one,
and the raven and nonblack marginals are \(1/100\) and \(4/5\) given
either alternative. For \(A\) (a black raven), \(P(A\mid H)=1/100\) and
\(P(A\mid\neg H)=1/200\), so \(R_L(A\Yright H)=2\). For \(B\) (a
nonblack nonraven), \(P(B\mid H)=4/5=160/200\) and
\(P(B\mid\neg H)=159/200\), so \(R_L(B\Yright H)=160/159\), which exceeds
the neutral value one by \(1/159\). By the first part, the same values
hold with \(G\) in place of \(H\).
\end{proof}
\end{annexitem}

\begin{annexitem}{annex:p15-w02c-p15-027-linda}{thm:p15-w02c-p15-027-linda}{The conjunction law, the representativeness ranking and the model}
\begin{proof}
\emph{Item 1.} \(BF\subseteq B\), so \(P(BF\cap E)\leq P(B\cap E)\);
divide by \(P(E)>0\).

\emph{Item 2.} All three events lie strictly between probability zero and
one, so the comparison is in the domain of POP-D
(Section~\ref{sec:p15-r43-axioms}). Applied with fixed target \(E\) and
sources \(BF\) and \(B\), POP-D states exactly that
\(R(BF\Yright E)_P\bowtie R(B\Yright E)_P\) holds when
\(P(E\mid BF)\bowtie P(E\mid B)\). For Q,
\(R_Q(X\Yright E)_P=P(E\mid X)/P(E)\) for \(X=BF,B\), and dividing by the
common positive \(P(E)\) preserves each of \(<,=,>\).

\emph{Item 3.} The event \(B\) is the disjoint union of \(BF\) and
\(B\neg F\). Put \(w=P(BF)/P(B)\), which lies in \((0,1)\). Then
\[
\begin{aligned}
 P(E\mid B)&=\frac{P(E\cap BF)+P(E\cap B\neg F)}{P(B)}\\
 &=w\,P(E\mid BF)+(1-w)\,P(E\mid B\neg F),
\end{aligned}
\]
so that
\[
 P(E\mid BF)-P(E\mid B)
 =(1-w)\bigl(P(E\mid BF)-P(E\mid B\neg F)\bigr),
\]
with \(1-w>0\).

\emph{Item 4.} The joint masses with \(E\) are \(P(E)\) times the
conditional masses: \(1/100\), \(1/200\), \(7/100\), \(3/200\). Subtracting
from the prior masses gives the masses with \(\neg E\): \(1/100\),
\(7/40\), \(41/100\), \(61/200\), all nonnegative, so the model is a
probability measure on eight atoms. Then \(P(BF)=1/50\), \(P(B)=1/5\),
\(P(F)=1/50+12/25=1/2\), and
\[
\begin{aligned}
 P(E\mid BF)&=\frac{1/100}{1/50}=\frac12,\qquad
 P(E\mid F)=\frac{1/100+7/100}{1/2}=\frac{4}{25},\\
 P(E\mid B)&=\frac{1/100+1/200}{1/5}=\frac{3}{40}.
\end{aligned}
\]
Subtracting \(P(E)=1/10\) gives the D values \(2/5\), \(3/50\) and
\(-1/40\); the second is positive and the third negative, so \(F\) is
positive and \(B\) negative to \(E\). Finally
\(P(BF\mid E)=1/10\) and \(P(B\mid E)=1/10+1/20=3/20\). The hypotheses of
items~2 and~3 hold, and \(P(E\mid B\neg F)=(1/200)/(9/50)=1/36<1/2\), in
agreement with item~3.
\end{proof}
\end{annexitem}

\begin{annexitem}{annex:p15-w02c-p15-027-base-rate}{thm:p15-w02c-p15-027-base-rate}{Posterior odds}
\begin{proof}
Write \(h=P(H)\in(0,1)\), \(a=P(E\mid H)\), \(b=P(E\mid\neg H)\) and
\(k=O(H)=h/(1-h)\in(0,\infty)\). Then \(P(E)=ha+(1-h)b>0\), so \(a\) and
\(b\) are not both zero, and \(P(H\mid E)=ha/P(E)\),
\(P(\neg H\mid E)=(1-h)b/P(E)\).

\emph{Item 1.} If \(b>0\), then
\(O(H\mid E)=ha/((1-h)b)=k\cdot a/b=k\,R_L(E\Yright H)_P\) by
Definition~\ref{def:p15-r43-measures}. If \(b=0\), then \(a>0\),
\(P(\neg H\mid E)=0\), so \(O(H\mid E)=+\infty\); and
\(R_L(E\Yright H)_P=+\infty\) by the same definition, so
\(k\cdot R_L=+\infty\) as \(k>0\).

\emph{Item 2.} Let \(r=R_L(E\Yright H)_P\). If \(0<r<\infty\), then
\(O(H\mid E)=kr\) equals \(r\) exactly when \(k=1\). If \(r=0\), then
\(O(H\mid E)=k\cdot0=0=r\). If \(r=+\infty\), then
\(O(H\mid E)=+\infty=r\). So \(r=O(H\mid E)\) exactly when \(k=1\) or
\(r\in\{0,+\infty\}\); and \(k=1\) exactly when \(h=1-h\), that is
\(h=1/2\). For the probability form, the map \(\tau(x)=x/(1+x)\) on
\([0,\infty)\), extended by \(\tau(+\infty)=1\), is strictly increasing
and so injective, and \(P(H\mid E)=\tau(O(H\mid E))\): for finite
\(O(H\mid E)=p/(1-p)\) with \(p=P(H\mid E)<1\), \(\phi\) returns \(p\),
and \(O(H\mid E)=+\infty\) exactly when \(p=1\). Hence
\(\tau(r)=P(H\mid E)\) exactly when \(r=O(H\mid E)\).

\emph{Item 3.} If \(0<r<\infty\), then \(O(H\mid E)/r=k\), which
exceeds one exactly when \(h>1-h\) and is below one exactly when
\(h<1-h\). The value \(r=a/b\) is determined by \(a\) and \(b\), and
\(h\) does not enter it.

\emph{Item 4.} The masses are nonnegative and sum to
\((9+81+3+7)/100=1\). \(P(H)=9/100+81/100=9/10\), so \(k=9\);
\(P(E\mid H)=(9/100)/(9/10)=1/10\) and
\(P(E\mid\neg H)=(3/100)/(1/10)=3/10\), so \(r=1/3<1\), and \(E\) is
negative to \(H\) because L preserves polarity with neutral value one
(Theorem~\ref{thm:p15-r43-measures}). By item~1,
\(O(H\mid E)=9\cdot1/3=3\), so \(P(H\mid E)=3/4\); directly,
\(P(H\mid E)=(9/100)/(12/100)=3/4\). Finally
\(\tau(1/3)=(1/3)/(4/3)=1/4\). Since \(3/4>1/2\), the posterior makes
\(H\) more probable than \(\neg H\), and since \(1/4<1/2\) the reply
makes \(\neg H\) more probable than \(H\).
\end{proof}
\end{annexitem}

\begin{annexitem}{annex:p15-w02c-p15-027-mixture}{thm:p15-w02c-p15-027-mixture}{Theorem~10 cases}
\begin{proof}
Throughout, \(P(AB)=0\), so for every event \(X\)
\[
 P((A\vee B)X)=P(AX)+P(BX)-P(ABX)=P(AX)+P(BX).
 \tag{\(\ast\)}
\]
The neutral values are \(\iota_D=0\) and \(\iota_Q=\iota_L=1\). If all
three values in \((M_\lambda)\) equal \(\iota_T\), the right side is
\(\lambda\iota_T+(1-\lambda)\iota_T=\iota_T\) for every
\(\lambda\in[0,1]\), so \((M_\lambda)\) holds for every such \(\lambda\).

\emph{Item 1.} Let \(P(C)\in\{0,1\}\) and let \(E\) be any event. If
\(P(E)=0\), the completions of Definition~\ref{def:p15-r43-measures}
give \(R_D=0\), \(R_Q=1\) and, as the target is degenerate, \(R_L=1\).
If \(P(E)>0\): when \(P(C)=0\), \(P(EC)\leq P(C)=0\), so
\(P(C\mid E)=0=P(C)\); when \(P(C)=1\), \(P(E\neg C)\leq P(\neg C)=0\), so
\(P(EC)=P(E)\) and \(P(C\mid E)=1=P(C)\). Hence \(R_D=P(C\mid E)-P(C)=0\).
For Q, the completion gives one when \(P(C)=0\), and
\(R_Q=P(C\mid E)/P(C)=1\) when \(P(C)=1\). For L the target is certain or
impossible, so the completion gives one. Thus every value is
\(\iota_T\), and the first paragraph gives \((M_\lambda)\) for every
\(\lambda\). Extreme positivity of \(A\) to \(C\) requires \(P(C)<1\) and
\(P(C\mid A)=1\), which fails when \(P(C)=1\), and when \(P(C)=0\)
because then \(P(C\mid A)=0\) if defined; extreme negativity is extreme
positivity to \(\neg C\), which has probability one or zero, so it fails
in the same way. If \(P(A)+P(B)>0\), the D and Q weight
\(P(A)/(P(A)+P(B))\) lies in \([0,1]\). When \(P(C)=0\),
\(P(\neg C)=1\) and \(P(X\mid\neg C)=P(X\neg C)=P(X)\) for every \(X\),
since \(P(XC)=0\); so the L weight equals \(P(A)/(P(A)+P(B))\). When
\(P(C)=1\), \(P(\neg C)=0\) and \(P(\cdot\mid\neg C)\) is undefined.

\emph{Item 2.} Let \(P(A)=P(B)=0\). Then \(P(A\vee B)=0\) by
\((\ast)\) with \(X\) the sure event. For any source \(E\) of probability
zero, \(R_D=0\) and \(R_Q=1\) by completion; for L, if the target is
degenerate the completion gives one, and otherwise
\(P(E\mid C)=P(E\mid\neg C)=0\), so both likelihoods vanish and the
completion gives one. Applying this to \(E=A\), \(B\) and \(A\vee B\),
all three values are \(\iota_T\), and \((M_\lambda)\) holds for every
\(\lambda\). The D and Q weight is \(0/(0+0)\). If \(P(C)<1\),
\(P(A\mid\neg C)=P(B\mid\neg C)=0\) and the L weight is \(0/0\); if
\(P(C)=1\) it is undefined as in item~1.

\emph{Item 3.} Let \(P(B)=0<P(A)\). Every event \(X\) has
\(P(BX)=0\), so by \((\ast)\) \(P((A\vee B)X)=P(AX)\) for \(X=C\) and
\(X=\neg C\), and also \(P(A\vee B)=P(A)\). Hence the joint tables
\(\pi_P(A\vee B,C)\) and \(\pi_P(A,C)\) coincide: their first two cells
are equal by these identities, and the last two are \(P(C)\) and
\(P(\neg C)\) minus the first two. Each of D, Q and L is a function of
the joint table, by Definition~\ref{def:p15-r43-measures}, so
\(R_T(A\vee B\Yright C)_P=R_T(A\Yright C)_P\). By item~2's argument for a
null source, \(R_T(B\Yright C)_P=\iota_T\), which is finite. So with
\(\lambda=1\) the right side of \((M_\lambda)\) is
\(1\cdot R_T(A\Yright C)_P+0\cdot\iota_T=R_T(A\Yright C)_P\), including
when that value is \(+\infty\). The D and Q weight is
\(P(A)/(P(A)+0)=1\). If \(P(C)<1\), then \(P(B\mid\neg C)=0\) and the L
weight is \(P(A\mid\neg C)/P(A\mid\neg C)=1\) whenever
\(P(A\mid\neg C)>0\).

\emph{Item 4.} Let \(0<P(C)<1\) and \(P(A),P(B)>0\), and write
\(a_X=P(X\mid C)\) and \(b_X=P(X\mid\neg C)\) for \(X=A,B\). By
\((\ast)\), \(P(A\vee B\mid C)=a_A+a_B\) and
\(P(A\vee B\mid\neg C)=b_A+b_B\). For \(X=A,B\): \(X\) is extremely
positive to \(C\) exactly when \(b_X=0\). Indeed, if \(b_X=0\) then
\(P(X\neg C)=0\), so \(P(C\mid X)=1>P(C)\); conversely
\(P(C\mid X)=1\) gives \(P(X\neg C)=0\). When \(b_X=0\),
\(a_X=P(X)/P(C)>0\), because \(P(XC)=P(X)\).

(a) If neither disjunct is extremely positive, \(b_A,b_B>0\), so
\(R_L(A\Yright C)_P=a_A/b_A\), \(R_L(B\Yright C)_P=a_B/b_B\) and
\(R_L(A\vee B\Yright C)_P=(a_A+a_B)/(b_A+b_B)\), all finite. With
\(\lambda=b_A/(b_A+b_B)\in[0,1]\), which is the L weight of part~(b) and
is defined because its denominator is positive,
\[
 \lambda\frac{a_A}{b_A}+(1-\lambda)\frac{a_B}{b_B}
  =\frac{a_A}{b_A+b_B}+\frac{a_B}{b_A+b_B}
  =R_L(A\vee B\Yright C)_P.
\]
Nothing here excludes \(a_A=0\) or \(a_B=0\), which is extreme negativity
of that disjunct to \(C\).

(b) Let \(b_A=0<b_B\). Then \(a_A>0\), so \(R_L(A\Yright C)_P=+\infty\),
while \(R_L(B\Yright C)_P=a_B/b_B\) and
\(R_L(A\vee B\Yright C)_P=(a_A+a_B)/b_B\) are finite and the second is
larger, as \(a_A>0\). For \(\lambda=0\) the right side of
\((M_\lambda)\) is \(R_L(B\Yright C)_P\), which differs from the left
side. For \(\lambda\in(0,1]\) it is \(+\infty\) plus the finite
nonnegative number \((1-\lambda)R_L(B\Yright C)_P\), that is,
\(+\infty\), which also differs from the finite left side.

(c) Let \(b_A=b_B=0\). Then \(a_A,a_B>0\), so both disjuncts and, as
\(b_A+b_B=0<a_A+a_B\), also \(A\vee B\) have L value \(+\infty\). For
\(\lambda\in(0,1]\) the right side is \(+\infty\) plus a nonnegative
term, and for \(\lambda=0\) it is \(0+R_L(B\Yright C)_P=+\infty\); so
\((M_\lambda)\) holds for every \(\lambda\). The L weight is
\(0/(0+0)\).
\end{proof}
\end{annexitem}

\begin{annexitem}{annex:p15-w02c-p15-027-lemma3}{thm:p15-w02c-p15-027-lemma3}{Lemma~3 cases}
\begin{proof}
\emph{Item 1.} If \(P(H)\in\{0,1\}\), the target is certain or
impossible, so the completion of Definition~\ref{def:p15-r43-measures}
gives \(R_L=1\) for every source, including \(E_1\), \(E_2\) and
\(E_1E_2\), and \(R_{LL}=\ln1=0\). Then \(1=1\cdot1\) and \(0=0+0\).

\emph{Setting.} Let \(0<h=P(H)<1\) and \((E_1\perp E_2\mid\pm H)_P\), and
put \(a_i=P(E_i\mid H)\) and \(b_i=P(E_i\mid\neg H)\). The hypothesis
gives \(P(E_1E_2\mid H)=a_1a_2\) and \(P(E_1E_2\mid\neg H)=b_1b_2\), and
total probability gives
\[
 P(E_i)=ha_i+(1-h)b_i,\qquad P(E_1E_2)=ha_1a_2+(1-h)b_1b_2 .
\]
As \(h>0\) and \(1-h>0\), an event with likelihoods \((a,b)\) is null
exactly when \(a=b=0\). For a source with likelihoods \((a,b)\),
Definition~\ref{def:p15-r43-measures} gives L value \(a/b\) if \(b>0\),
\(+\infty\) if \(b=0<a\), and \(1\) if \(a=b=0\); the LL value is the
logarithm of the L value, with \(\ln0=-\infty\) and
\(\ln(+\infty)=+\infty\).
\emph{Fact F.} The L value is \(+\infty\) exactly when the source is
extremely positive to \(H\), and \(0\) exactly when it is extremely
negative. Indeed \(b=0<a\) means \(P(E\neg H)=0<P(E)\), that is,
\(P(E)>0\) and \(P(H\mid E)=1\), which with \(P(H)<1\) is extreme
positivity (Definition~\ref{def:p15-r43-extreme}); and \(a=0<b\) means
\(P(E)>0\) and \(P(\neg H\mid E)=1\), which with \(P(\neg H)<1\) is
extreme positivity to \(\neg H\), that is, extreme negativity.

\emph{Item 2.} Let \(P(E_1E_2)>0\), so \(a_1a_2>0\) or \(b_1b_2>0\). If
\(l_1=+\infty\) and \(l_2=0\), then \(b_1=0\) and \(a_2=0\), so
\(a_1a_2=b_1b_2=0\), a contradiction; by symmetry
\(\{l_1,l_2\}\neq\{0,+\infty\}\). As \(m_i=\ln l_i\), also
\(\{m_1,m_2\}\neq\{-\infty,+\infty\}\), so \(l_1l_2\) and \(m_1+m_2\) are
defined. Three cases exhaust the possibilities.
(i) \(b_1,b_2>0\). Then \(b_1b_2>0\), and
\(l_{12}=a_1a_2/(b_1b_2)=(a_1/b_1)(a_2/b_2)=l_1l_2\), a finite number. If
\(l_1,l_2>0\), taking logarithms gives \(m_{12}=m_1+m_2\). If some
\(l_i=0\), then \(l_{12}=0\), \(m_{12}=-\infty\), \(m_i=-\infty\) and the
other LL value is less than \(+\infty\), so \(m_1+m_2=-\infty=m_{12}\).
(ii) Exactly one of \(b_1,b_2\) is zero, say \(b_1=0<b_2\). Then
\(b_1b_2=0\), so \(a_1a_2>0\) and \(a_1,a_2>0\). Hence
\(l_{12}=+\infty\), \(l_1=+\infty\) and \(0<l_2=a_2/b_2<+\infty\), so
\(l_1l_2=+\infty=l_{12}\); and \(m_{12}=+\infty=m_1+m_2\), since \(m_2\)
is finite. The case \(b_2=0<b_1\) is the same with indices exchanged.
(iii) \(b_1=b_2=0\). Then \(a_1a_2>0\), and
\(l_1=l_2=l_{12}=+\infty\) and \(m_1=m_2=m_{12}=+\infty\), so both
identities hold.
For the particular statements: if a source is extremely negative, its
L value is \(0\) by Fact~F, so its \(b_i>0\) and \(a_i=0\); then
\(a_1a_2=0\), so \(b_1b_2>0\), case~(i) applies and \(l_{12}=0\). If a
source is extremely positive and the other is not extremely negative,
cases~(ii) or~(iii) apply and \(l_{12}=+\infty\).

\emph{Item 3.} Let \(P(E_1E_2)=0\). Then \(a_1a_2=b_1b_2=0\), both
likelihoods of \(E_1E_2\) vanish, and the completion gives \(l_{12}=1\)
and \(m_{12}=0\). The two cases below are exclusive and exhaustive.
(a) Let \(P(E_1)=0\), so \(a_1=b_1=0\), \(l_1=1\) and \(m_1=0\). The
identities become \(1=l_2\) and \(0=m_2\), which are equivalent. Now
\(l_2=1\) exactly when \(a_2=b_2\): if \(b_2>0\) this is \(a_2/b_2=1\),
if \(b_2=0<a_2\) the value is \(+\infty\), and if \(a_2=b_2=0\) it is
\(1\). Since
\(P(E_2H)-P(E_2)P(H)=ha_2-h\bigl(ha_2+(1-h)b_2\bigr)=h(1-h)(a_2-b_2)\),
\(a_2=b_2\) says exactly that \(E_2\) is irrelevant to \(H\), which
includes the case that \(E_2\) is null. The case \(P(E_2)=0\) is the same
with indices exchanged.
(b) Let \(P(E_1),P(E_2)>0\), so \((a_i,b_i)\neq(0,0)\). If \(a_1=0\),
then \(b_1>0\); from \(b_1b_2=0\), \(b_2=0\), so \(a_2>0\); hence
\(l_1=0\) and \(l_2=+\infty\). If \(a_1>0\), then \(a_1a_2=0\) gives
\(a_2=0\), so \(b_2>0\), and \(b_1b_2=0\) gives \(b_1=0\); hence
\(l_1=+\infty\) and \(l_2=0\). By Fact~F one source is extremely
positive and the other extremely negative; the LL values are
\(+\infty\) and \(-\infty\), and \(0\cdot(+\infty)\) and
\((+\infty)+(-\infty)\) are undefined.

\emph{Item 4.} If \(P(E)=0\), then \(P(EH)=0=P(E)P(H)\), so \(E\) is
irrelevant to \(H\) by Definition~\ref{def:p15-r43-polarity}; and
\(P(EE')\leq P(E)=0\), so \(EE'\) is irrelevant to \(H\) as well. Quality
Representation gives both the value \(\iota_T\). If \(\iota_T=1\), the
product identity reads \(1=1\cdot r\) with \(r=R_T(E'\Yright H)_P\); since
\(1\cdot r=r\) for every \(r\in\mathbb R^*\), with
\(1\cdot(\pm\infty)=\pm\infty\), it holds exactly when \(r=1=\iota_T\), that is, by Quality Representation,
exactly when \(E'\) is irrelevant to \(H\). If \(\iota_T=0\), the sum
identity reads \(0=0+r\), with \(0+(\pm\infty)=\pm\infty\), which holds
exactly when \(r=0=\iota_T\).

\emph{Item 5.} (i) With \(E_1\) null, \(P(E_1E_2\mid H)\leq P(E_1\mid H)=0
=a_1a_2\) and likewise given \(\neg H\), so the hypothesis holds. Realize
\(E_2\) and \(H\) by (R2) with \(h=1/2\), \(a=1/2\), \(b=1/4\), and let
\(E_1\) be empty. Then \(l_2=2\) and \(m_2=\ln2\), while \(l_1=1\) and
\(m_1=0\); item~3 gives \(l_{12}=1\) and \(m_{12}=0\), against
\(l_1l_2=2\) and \(m_1+m_2=\ln2\). (ii) On the four points,
\(P(H)=1/2\), \(a_1=1/2\), \(b_1=0\), \(a_2=0\) and \(b_2=1/2\). As
\(E_1E_2\) is empty, \(P(E_1E_2\mid H)=0=a_1a_2\) and
\(P(E_1E_2\mid\neg H)=0=b_1b_2\), so the hypothesis holds. The L values
are \(l_1=+\infty\) and \(l_2=0\), the LL values \(+\infty\) and
\(-\infty\), and item~3 gives \(l_{12}=1\) and \(m_{12}=0\).
\end{proof}
\end{annexitem}

\begin{annexitem}{annex:p15-w02c-p15-027-pooling}{thm:p15-w02c-p15-027-pooling}{Pooling}
\begin{proof}
Use the notation of the setting of the preceding proof. As
\(P(E_i)\geq P(E_1E_2)>0\),
\[
 q_i=\frac{ha_i}{P(E_i)},\qquad 1-q_i=\frac{(1-h)b_i}{P(E_i)} .
\]

\emph{Item 1.} Multiplying,
\[
 q_1q_2(1-h)=\frac{h^2(1-h)a_1a_2}{P(E_1)P(E_2)},\qquad
 (1-q_1)(1-q_2)h=\frac{h(1-h)^2b_1b_2}{P(E_1)P(E_2)} .
\]
Their sum is \(h(1-h)P(E_1E_2)/(P(E_1)P(E_2))>0\), which is the
denominator. Dividing the first quantity by the sum gives
\(ha_1a_2/P(E_1E_2)\). By the hypothesis
\(P(HE_1E_2)=hP(E_1E_2\mid H)=ha_1a_2\), so this quotient is
\(P(H\mid E_1E_2)=q_{12}\).

\emph{Item 2.} For a source with \(P(E)>0\) and \(0<h<1\), the
definitions give \(R_D=q-h\), \(R_Q=q/h\) and \(R_{QL}=\ln(q/h)\), which
is \(-\infty\) exactly when \(q=0\); so \(q=h+R_D=hR_Q=he^{R_{QL}}\). For
Z, \(R_Z=(q-h)/(1-h)>0\) when \(q>h\) and \(R_Z=(q-h)/h\leq0\) when
\(q\leq h\), so the sign of \(R_Z\) tells which formula applies, and
\(q=h+(1-h)R_Z\) or \(q=h+hR_Z\) accordingly. The two sources and the
conjunction all have positive probability, so their values are given by
these formulas and not by completions. Hence \(h\) and the constituent
values give \(q_1,q_2\); item~1 gives \(q_{12}\); and the same formula
applied to \(q_{12}\) gives \(R_T(E_1E_2\Yright H)_P\).

\emph{Item 3.} For \(M(h,a,b)\) take the sample space \(\{0,1\}^3\)
of triples \((x,y_1,y_2)\), with \(H=\{x=1\}\),
\(E_i=\{y_i=1\}\) and masses
\[
 \begin{cases}
  ha^{y_1+y_2}(1-a)^{2-y_1-y_2}, & x=1,\\
  (1-h)b^{y_1+y_2}(1-b)^{2-y_1-y_2}, & x=0.
 \end{cases}
\]
The masses are nonnegative and sum to \(h+(1-h)=1\); \(P(H)=h\),
\(P(E_i\mid H)=a\), \(P(E_1E_2\mid H)=a^2=a\cdot a\), and likewise with
\(b\) given \(\neg H\). So the hypotheses hold, with
\(P(E_1E_2)=ha^2+(1-h)b^2>0\) for the listed parameters, and
\[
 q_i=\frac{ha}{ha+(1-h)b},\qquad q_{12}=\frac{ha^2}{ha^2+(1-h)b^2}.
\]
The four models give
\[
\begin{array}{l|cc|c|cc}
 \text{model} & q_i & q_{12} & T & R_T(E_i\Yright H)_P
   & R_T(E_1E_2\Yright H)_P\\\hline
 M(1/2,3/4,1/4) & 3/4 & 9/10 & D & 1/4 & 2/5\\
 & & & Q & 3/2 & 9/5\\
 & & & Z & 1/2 & 4/5\\
 M(1/4,3/4,1/4) & 1/2 & 3/4 & D & 1/4 & 1/2\\
 M(1/3,1/2,1/4) & 1/2 & 2/3 & Q & 3/2 & 2\\
 M(1/4,5/6,1/6) & 5/8 & 25/28 & Z & 1/2 & 6/7
\end{array}
\]
In the order of the four models, the posteriors are computed as
\[
\begin{array}{l|cccc}
 q_i & (3/8)/(1/2) & (3/16)/(6/16) & (1/6)/(1/3) & (5/24)/(8/24)\\
 q_{12} & (9/32)/(10/32) & (9/64)/(12/64) & (1/12)/(1/8)
   & (25/144)/(28/144)
\end{array}
\]
Z uses its first
formula, since every posterior exceeds the prior: for instance
\((3/4-1/2)/(1/2)=1/2\), \((9/10-1/2)/(1/2)=4/5\),
\((5/8-1/4)/(3/4)=1/2\) and \((25/28-1/4)/(3/4)=6/7\). A function \(g_D\) as described would have \(g_D(1/4,1/4)\)
equal to both \(2/5\) and \(1/2\); likewise \(g_Q(3/2,3/2)\) would equal
\(9/5\) and \(2\), \(g_{QL}(\ln\tfrac32,\ln\tfrac32)\) would equal
\(\ln\tfrac95\) and \(\ln2\), and \(g_Z(1/2,1/2)\) would equal \(4/5\)
and \(6/7\). No such functions exist.

\emph{Item 4.} The L and LL statements are item~2 of
Theorem~\ref{thm:p15-w02c-p15-027-lemma3}, and the KO statement is
item~1 of Theorem~\ref{thm:p15-w02b-p15-024-ko-transform}. In
\(M(h,3/4,1/4)\) the constituent L values are \((3/4)/(1/4)=3\) and the
conjunction's is \((9/16)/(1/16)=9\) for every \(h\), while the D values
of the conjunction are \(2/5\) for \(h=1/2\) and \(1/2\) for \(h=1/4\).
\end{proof}
\end{annexitem}

\begin{annexitem}{annex:p15-w02c-p15-027-homomorphism}{thm:p15-w02c-p15-027-homomorphism}{A certain conjunct}
\begin{proof}
Let \(P(B)=1\). For any event \(X\), \(P(X\neg B)\leq P(\neg B)=0\), so
\(P(XB)=P(X)\). With \(X=H\), \(\Delta_P(B,H)=P(H)-1\cdot P(H)=0\), so
\(B\) is irrelevant to \(H\) and Quality Representation gives
\(R_T(B\Yright H)_P=\iota_T\). With \(X=AH\) and \(X=A\neg H\),
\(P(ABH)=P(AH)\) and \(P(AB\neg H)=P(A\neg H)\); subtracting these from
\(P(H)\) and \(P(\neg H)\) gives \(P(\neg(AB)H)=P(\neg AH)\) and
\(P(\neg(AB)\neg H)=P(\neg A\neg H)\). So
\(\pi_P(AB,H)=\pi_P(A,H)\), and Transparency gives
\(R_T(AB\Yright H)_P=R_T(A\Yright H)_P\). If \(\iota_T=0\), adding
\(R_T(B\Yright H)_P=0\) changes nothing, including at infinite values;
if \(\iota_T=1\) and \(R_T(A\Yright H)_P\) is finite, multiplying by
\(R_T(B\Yright H)_P=1\) changes nothing.
\end{proof}
\end{annexitem}
